\documentclass[12pt]{article}
\usepackage{amsfonts,amsthm,amsmath,amssymb,mathtools}
\usepackage{mathrsfs}
\usepackage{enumitem}
\usepackage{tikz-cd}
\usepackage{geometry}
\usepackage{mlmodern}

\geometry{letterpaper, portrait, margin=1in}

\setcounter{section}{0}

\renewcommand{\thesection}{\S\arabic{section}}
\renewcommand{\thesubsection}{\arabic{section}}
\newcommand{\dd}{\mathop{}\!\mathrm{d}}

\newtheorem{thm}{Theorem}
\newenvironment{theorem}[1]{
	\renewcommand{\thethm}{#1}
	\begin{thm}
		}{\end{thm}}

\newtheorem{lma}{Lemma}
\newenvironment{lemma}[1]{
	\renewcommand{\thelma}{#1}
	\begin{lma}
		}{\end{lma}}

\theoremstyle{definition}
\newtheorem{ex}{}
\newenvironment{exercise}[1]{
	\renewcommand{\theex}{#1}
	\begin{ex}
		}{\end{ex}}

\title{Homework 5}
\author{Connor Mooneyhan}
\date{October 2, 2025}

\begin{document}

\maketitle

\begin{ex}
	If $G$ is a finitely-generated abelian group, $G$ is isomorphic to a direct
	product of cyclic groups. However, the decomposition as a product of cyclic
	groups is not in general unique.
	\begin{enumerate}[label=(\alph*)]
		\item Find an element of order $30$ in the group $\mathbb{Z}/5\mathbb{Z} \times \mathbb{Z}/6\mathbb{Z}$.
		\item Use your answer to (a) to prove that $\mathbb{Z}/5\mathbb{Z} \times \mathbb{Z}/6\mathbb{Z} \cong \mathbb{Z}/30\mathbb{Z}$.
	\end{enumerate}
\end{ex}
\begin{proof}\
	\begin{enumerate}[label=(\alph*)]
		\item The element $(1, 1)$ has order 30.
		\item Since the order of $\mathbb{Z}/5 \mathbb{Z} \times \mathbb{Z}/6 \mathbb{Z}$ is 30, part (a) shows that it is cyclic since it is generated by one element.
		      Thus $\mathbb{Z}/5\mathbb{Z} \times \mathbb{Z}/6\mathbb{Z} \cong \mathbb{Z}/30\mathbb{Z}$ since cyclic groups of the same order are isomorphic.
	\end{enumerate}
\end{proof}

\begin{ex}
	Let $C$ be a non-singular cubic curve given by $y^{2} = f(x) = x^{3} + ax^{2} + bx + c$.
	\begin{enumerate}[label=(\alph*)]
		\item Prove that
		      \begin{equation*}
			      \frac{d^{2} y}{dx^{2}} = \frac{2f''(x) f(x) - f'(x)^{2}}{4yf(x)} = \frac{\psi_{3}(x)}{4yf(x)}.
		      \end{equation*}
		      Use this to deduce that a point $P = (x,y) \in C$ is a point of order $3$
		      if and only if $P \ne \mathcal{O}$ and $P$ is a point of inflection on the curve $C$.
		\item Suppose now that $a, b, c \in \mathbb{R}$. Prove that $\psi_{3}(x)$ has
		      exactly two real roots, say $\alpha_{1}$ and $\alpha_{2}$ with
		      $\alpha_{1} < \alpha_{2}$. Show that $f(\alpha_{1}) < 0$ and
		      $f(\alpha_{2}) > 0$. (From this it follows that $C[3](\mathbb{R}) \cong \mathbb{Z}/3\mathbb{Z}$.)
	\end{enumerate}
\end{ex}
\begin{proof}\
	\begin{enumerate}[label=(\alph*)]
		\item Indeed, \begin{align*}
			      f(x)   & = x^3 + ax^2 + bx + c = y^2               \\
			      f'(x)  & = 3x^2 + 2ax + b = 2y \frac{\dd y}{\dd x} \\
			      f''(x) & = 6x + 2a
		      \end{align*}
		      and \begin{align*}
			      y^2                              & = x^3 + ax^2 + bx + c                                         \\
			      \implies 2y \frac{\dd y}{\dd x}  & = 3x^2 + 2ax + b                                              \\
			      \implies \frac{\dd y}{\dd x}     & = \frac{3x^2 + 2ax + b}{2y}                                   \\
			      \implies \frac{\dd^2 y}{\dd x^2} & = \frac{(6x+2a)2y - 2\frac{\dd y}{\dd x}(3x^2 + 2ax+b)}{4y^2} \\
			                                       & = \frac{f''(x)2y - \frac{f'(x)}{y}f'(x)}{4f(x)}               \\
			                                       & = \frac{f''(x)2y^2 - f'(x)^2}{4yf(x)}                         \\
			                                       & = \frac{2f''(x)f(x) - f'(x)^2}{4yf(x)}                        \\
                                             & = \frac{\psi_3(x)}{4yf(x)}.
		      \end{align*}
          Certainly, $P$ is of order 1 if and only if $P = \mathcal{O}$ since $\mathcal{O}$ is the identity, so $P \neq \mathcal{O}$ if $P$ is of order 3.
          By definition, points of inflection are points where the second derivative is equal to 0, and the expression above is equal to 0 if and only if $\psi_3(x) = 0$ and $y \neq 0$ (since $y = 0$ if and only if $f(x) = 0$).
          Indeed, if $y$ were 0, then $P$ would be of order 2, not 3, so we can omit the restriction on $y$ and determine that the points of order 3 are exactly the points of inflection which are not $\mathcal{O}$.
        \item 
	\end{enumerate}
\end{proof}

\begin{ex}
	$*$ Construct yourself an elliptic curve with as high rank as you can.
	For MTH 346 students, full credit will be awarded for a curve of rank $\geq 4$.
	For MTH 646 students, full credit will be awarded for a curve of rank $\geq 5$.

	You should use Magma to test the rank of your elliptic curve.

	When I say construct such a curve yourself, I mean you should *not* just
	look up curves in databases of curves with high rank. You should use some
	method to search. Just testing all curves in a box (say those of the form
	$y^{2} = x^{3} + Ax + B$ where $A, B \in \mathbb{Z}$ and $0 \leq A,B \leq 100$) will likely work for finding a rank $4$ curve. For rank $\geq 5$ you need a way to decrease the number of curves
	you are testing, or find a systematic way to construct curves with lots of rational points on them.
\end{ex}

\end{document}
